% Options for packages loaded elsewhere
\PassOptionsToPackage{unicode}{hyperref}
\PassOptionsToPackage{hyphens}{url}
\documentclass[
  12pt,
  a4paper,
]{article}
\usepackage{xcolor}
\usepackage[top=20mm,bottom=20mm,left=25mm,right=25mm]{geometry}
\usepackage{amsmath,amssymb}
\setcounter{secnumdepth}{-\maxdimen} % remove section numbering
\usepackage{iftex}
\ifPDFTeX
  \usepackage[T1]{fontenc}
  \usepackage[utf8]{inputenc}
  \usepackage{textcomp} % provide euro and other symbols
\else % if luatex or xetex
  \usepackage{unicode-math} % this also loads fontspec
  \defaultfontfeatures{Scale=MatchLowercase}
  \defaultfontfeatures[\rmfamily]{Ligatures=TeX,Scale=1}
\fi
\usepackage{lmodern}
\ifPDFTeX\else
  % xetex/luatex font selection
  \setmainfont[]{Times New Roman}
\fi
% Use upquote if available, for straight quotes in verbatim environments
\IfFileExists{upquote.sty}{\usepackage{upquote}}{}
\IfFileExists{microtype.sty}{% use microtype if available
  \usepackage[]{microtype}
  \UseMicrotypeSet[protrusion]{basicmath} % disable protrusion for tt fonts
}{}
\makeatletter
\@ifundefined{KOMAClassName}{% if non-KOMA class
  \IfFileExists{parskip.sty}{%
    \usepackage{parskip}
  }{% else
    \setlength{\parindent}{0pt}
    \setlength{\parskip}{6pt plus 2pt minus 1pt}}
}{% if KOMA class
  \KOMAoptions{parskip=half}}
\makeatother
\setlength{\emergencystretch}{3em} % prevent overfull lines
\providecommand{\tightlist}{%
  \setlength{\itemsep}{0pt}\setlength{\parskip}{0pt}}
\usepackage{setspace}
\usepackage{graphicx}
\usepackage{booktabs,array}
\setstretch{1.15}
\setlength{\parindent}{0pt}
\setlength{\parskip}{4pt plus 1pt}
\usepackage{titlesec}
\titleformat{\section}{\large\bfseries}{}{0pt}{}
\titleformat{\subsection}{\normalsize\bfseries}{}{0pt}{}
\titleformat{\subsubsection}{\normalsize\bfseries}{}{0pt}{}
\titlespacing*{\section}{0pt}{0.25\baselineskip}{0.2\baselineskip}
\titlespacing*{\subsection}{0pt}{0pt}{0pt}
\titlespacing*{\subsubsection}{0pt}{0pt}{0pt}
\usepackage{caption}
\captionsetup{font=small,labelfont=bf,width=0.95\linewidth,justification=centering,skip=3pt}
\usepackage{enumitem}
\setlist{nosep,topsep=2pt}
\setlength{\textfloatsep}{8pt plus 2pt minus 2pt}
\setlength{\floatsep}{6pt plus 2pt minus 2pt}
\setlength{\intextsep}{6pt plus 2pt minus 2pt}
\renewcommand{\topfraction}{0.9}
\renewcommand{\bottomfraction}{0.85}
\renewcommand{\textfraction}{0.1}
\renewcommand{\floatpagefraction}{0.75}
\setcounter{topnumber}{3}
\setcounter{bottomnumber}{3}
\setcounter{totalnumber}{5}
\newcommand{\proofslot}[2]{\par\noindent\fbox{\parbox[c][#1][c]{\dimexpr\linewidth-2\fboxsep-2\fboxrule\relax}{\centering #2}}\par}
\newcommand{\proofpage}{\noindent\textbf{Mesh-sensitivity results: one-page allowance}\par\medskip}
\usepackage{bookmark}
\IfFileExists{xurl.sty}{\usepackage{xurl}}{} % add URL line breaks if available
\urlstyle{same}
\hypersetup{
  hidelinks,
  pdfcreator={LaTeX via pandoc}}

\author{}
\date{}

\begin{document}

\section{1. Introduction}\label{introduction}

\subsection{1.1. Engineering problem}\label{engineering-problem}

Geothermal steam--water separators supply dry steam to turbines. Liquid
carryover can cause erosion and transport minerals that form deposits.
Design must therefore consider separation efficiency, pressure drop and
operating pressure.
(\href{https://doi.org/10.1016/j.geothermics.2014.05.009}{Zarrouk \&
Purnanto, 2015})

This study develops an inherited computational fluid dynamics (CFD)
model of a vertical bottom-outlet cyclone with a spiral inlet. It
follows liquid from inlet placement through wall transport and
collection.

\subsection{1.2. Problem in the inherited
model}\label{problem-in-the-inherited-model}

The simplified geometry omits the lower brine-discharge path. A closed
bottom gives liquid no intended exit; opening the whole lower boundary
can also let steam escape there. The model needs a collection treatment
that preserves the steam route.

Liquid reaching the wall is only one stage of separation. A lower liquid
flow through the steam outlet can accompany greater collection,
continued storage or transfer into a wall film. The model must therefore
follow liquid destinations alongside outlet flow to assess what each
change achieves.

A virtual collector, termed an absorber, and reduced-feed startup
enabled roughness and Eulerian Wall Film (EWF) studies after
unsuccessful full-geometry drainage trials.

\subsection{1.3. Research aim and
questions}\label{research-aim-and-questions}

The aim is to develop and assess the separator model's liquid-transport
and collection paths. Predicting separation efficiency and pressure drop
is the longer-term purpose. The present work addresses four questions:

\begin{enumerate}
\def\labelenumi{\arabic{enumi}.}
\tightlist
\item
  How do mesh representation and inlet phase placement affect liquid
  distribution and retention near the wall?
\item
  How can numerical collection and staged startup provide starting
  fields for wall studies?
\item
  How do wall roughness and wall-film modelling affect liquid transport,
  storage, release and collection?
\item
  What numerical and liquid-accounting evidence supports these
  predictions, and what remains necessary for field validation?
\end{enumerate}

\subsection{1.4. Scope and contribution}\label{scope-and-contribution}

The study mainly uses the simplified primary separator. Mesh and inlet
comparisons lead into collector development, startup and wall-treatment
studies. Later cases examine film development to 500 ms, release
mechanisms and direct drainage. Downstream equipment, mineral chemistry
and optimum operating pressure are outside the simulated scope.

The contribution combines checked collection and startup procedures with
evidence of how model choices affect liquid delivery, storage and
removal. The comparisons assess model development rather than validated
separator performance; refinement of the combined model remains
incomplete. Appendix A.2 defines the controlled reconstructions.

Figure 1 maps the liquid destinations examined across separate
configurations. Transfer into film or droplets changes the liquid
representation within the vessel; collector and drain treatments remove
it from the model. Released-droplet destinations require tracking.
During frozen-bulk film tests, reciprocal depletion of the bulk liquid
is not advanced.

\begin{figure}[htbp]
\centering
\includegraphics[width=16cm]{\detokenize{/Users/shuheiyokkaichi/Developer/P4P_knowledgeWiki/Project/Final research report/Shuhei Report/figures/report-liquid-destinations-overview.png}}
\caption*{Figure 1. Liquid destinations examined in separate model configurations.
Wall capture can lead to storage, release or removal. Numerical sinks
and native film-edge outflow are distinct; this is a conceptual map, not
a measured balance.}
\end{figure}

Section 3 describes the model and comparisons. Section 4 asks how liquid
reaches the wall, where it can go and what wall treatment changes.
Sections 5--6 connect these findings to collection and the evidence
needed for performance prediction.

\section{2. Literature review}\label{literature-review}

\subsection{2.1. How a geothermal steam--water separator
works}\label{how-a-geothermal-steamwater-separator-works}

A vertical bottom-outlet cyclone (BOC) uses a tangential or spiral inlet
to rotate the incoming steam and brine. Liquid moves towards the wall
while steam occupies the central region. An internal pipe carries steam
out through the bottom; brine leaves separately. \emph{Bottom outlet}
refers to the steam-pipe arrangement.
(\href{https://doi.org/10.1016/j.geothermics.2014.05.009}{Zarrouk \&
Purnanto, 2015})

Liquid at the wall can form a film that drains towards the brine
discharge or releases droplets back into the steam (Figure 2). The
schematic also distinguishes brine discharge from downstream liquid
collection and condensation.

\begin{figure}[htbp]
\centering
\includegraphics[width=7cm]{\detokenize{/Users/shuheiyokkaichi/Developer/P4P_knowledgeWiki/Project/Final research report/Shuhei Report/figures/literature/rizaldy-2016-figure-5-operation.png}}
\caption*{Figure 2. Conceptual separator liquid paths: wall film, carryover, brine
discharge (\(\dot m_l\)) and downstream collection (\(\dot m_b\)).
Source:
\href{https://www.researchgate.net/publication/310818568_LIQUID_CARRYOVER_IN_GEOTHERMAL_STEAM-WATER_SEPARATORS}{Rizaldy
et al.~(2016)}, fig.~5.}
\end{figure}

\subsection{2.2. Internal flow physics and performance
measures}\label{internal-flow-physics-and-performance-measures}

Rotation produces a low-pressure core, with inward acceleration scaling
as \(u_\theta^2/r\) for tangential velocity \(u_\theta\) and radius
\(r\). Density-dependent responses to pressure and drag allow outward
liquid transport. Fine droplets can follow steam; larger droplets reach
collection surfaces more readily. Uncertain inlet sizes and continued
droplet generation complicate prediction.
(\href{https://publications.mygeoenergynow.org/grc/1028587.pdf}{Pointon
et al., 2009})

Higher velocity can strengthen outward transport but also increase
entrainment by steam shear. Empirical data show rapid wetness increases
at different velocities for different separators (Figure 3). Thus,
greater swirl alone does not establish better separation.
(\href{https://doi.org/10.1016/j.geothermics.2014.05.009}{Zarrouk \&
Purnanto, 2015})

\begin{figure}[htbp]
\centering
\includegraphics[width=10cm]{\detokenize{/Users/shuheiyokkaichi/Developer/P4P_knowledgeWiki/Project/Final research report/Shuhei Report/figures/literature/zarrouk-purnanto-2015-figure-12-breakdown.png}}
\caption*{Figure 3. Outlet wetness rises at different inlet velocities across
three datasets. Right axis: steam quality. Source:
\href{https://doi.org/10.1016/j.geothermics.2014.05.009}{Zarrouk and
Purnanto (2015)}, fig.~12, using Bangma (1960) and Lazalde-Crabtree
(1984) data.}
\end{figure}

Outlet steam quality is

\[
x_{\mathrm{out}}=
\frac{\dot m_{v,\mathrm{out}}}
{\dot m_{v,\mathrm{out}}+\dot m_{l,\mathrm{out}}},
\]

where the terms are outward vapour and liquid mass-flow rates through
the steam outlet. Purnanto et al.~(2013) used this measure. It differs
from liquid-capture fraction, which compares collected liquid with inlet
liquid; their denominators are not interchangeable.
(\href{https://www.researchgate.net/publication/269519626_CFD_MODELLING_OF_TWO-PHASE_FLOW_INSIDE_GEOTHERMAL_STEAM-WATER_SEPARATORS}{Purnanto
et al., 2013})

Steam purity concerns contamination. Chemical carryover estimates are
affected by condensation and drain-pot collection, so field comparisons
need matching definitions and sampling locations.
(\href{https://www.researchgate.net/publication/310818568_LIQUID_CARRYOVER_IN_GEOTHERMAL_STEAM-WATER_SEPARATORS}{Rizaldy
et al., 2016})

Internal pressure drop requires defined inlet/outlet locations and a
stated static or total pressure measure. It differs from selecting
optimum operating pressure, which also involves flashing and plant
conditions. The present no-flashing model addresses internal flow at
prescribed phase conditions.
(\href{https://publications.mygeoenergynow.org/grc/1028587.pdf}{Pointon
et al., 2009};
\href{https://doi.org/10.1016/j.geothermics.2014.05.009}{Zarrouk \&
Purnanto, 2015})

\subsection{2.3. Established design methods and geothermal
CFD}\label{established-design-methods-and-geothermal-cfd}

Empirical geothermal design methods provide reference points for CFD.
However, CFD's additional flow detail must be assessed against its
assumptions and comparison evidence.
(\href{https://doi.org/10.1016/j.geothermics.2014.05.009}{Zarrouk \&
Purnanto, 2015})

Pointon et al.~(2009) used Fluent's Discrete Phase Model (DPM) with
assumed droplet feed and wall/baffle collection. Predicted efficiencies
were 99.96\% for scrolled and 99.93\% for tangential entry; an empirical
design method gave 99.955\%. Scrolled entry also gave higher mean
tangential velocities, 34.1 rather than 32.5 m/s (Figure 4). Empirical
agreement did not independently test the feed or collection assumptions.
(\href{https://publications.mygeoenergynow.org/grc/1028587.pdf}{Pointon
et al., 2009})

\begin{figure}[htbp]
\centering
\emph{(a) Inlet configurations.}
\par
\includegraphics[width=10cm]{\detokenize{/Users/shuheiyokkaichi/Developer/P4P_knowledgeWiki/Project/Final research report/Shuhei Report/figures/literature/pointon-2009-figure-5-inlets.png}}
\par
\emph{(b) Calculated tangential-velocity distributions.}
\par
\includegraphics[width=10cm]{\detokenize{/Users/shuheiyokkaichi/Developer/P4P_knowledgeWiki/Project/Final research report/Shuhei Report/figures/literature/pointon-2009-figure-7-velocity.png}}
\caption*{Figure 4. Scrolled entry gives higher calculated tangential velocity.
(a) Geometries; (b) velocity fields. Scrolled: left; tangential: right.
Source:
\href{https://publications.mygeoenergynow.org/grc/1028587.pdf}{Pointon
et al.~(2009)}, figs. 5 and 7.}
\end{figure}

Purnanto et al.~(2013) compared three BOC configurations using a steady
Mixture model and DPM outlet-quality calculation. Changed vessel
proportions prevented isolation of inlet shape. Counting incomplete
droplet tracks as separated improved empirical agreement, but the
authors called that assumption non-rigorous. The inherited modelling
basis therefore requires explicit treatment of unresolved droplet fate.
(\href{https://www.researchgate.net/publication/269519626_CFD_MODELLING_OF_TWO-PHASE_FLOW_INSIDE_GEOTHERMAL_STEAM-WATER_SEPARATORS}{Purnanto
et al., 2013})

Both studies used renormalisation-group (RNG) k--ε turbulence modelling.
Pointon's swirl modification and lower-cost rationale offer a starting
point, not a detailed turbulence-model comparison. Accuracy in the
present swirling and near-wall flow remains to be assessed.
(\href{https://publications.mygeoenergynow.org/grc/1028587.pdf}{Pointon
et al., 2009};
\href{https://www.researchgate.net/publication/269519626_CFD_MODELLING_OF_TWO-PHASE_FLOW_INSIDE_GEOTHERMAL_STEAM-WATER_SEPARATORS}{Purnanto
et al., 2013})

Appendix B.2 compares the methods and assumptions carried into this
study.

\subsection{2.4. Wall--liquid transport and
carryover}\label{wallliquid-transport-and-carryover}

Film motion depends on gravity, pressure and shear. Deposition adds
liquid; drainage removes it; entrainment returns droplets to the steam.
More wall deposition can therefore increase storage or later release
without increasing external collection.

Rizaldy et al.~(2016) predict entrainment sensitivity to loading, film
thickness and velocity. Figure 5 shows undercut and roll-wave release,
associated with lower and higher film Reynolds numbers, respectively.
Film Reynolds number compares inertia with viscosity using liquid
velocity and film thickness. Missing brine-flow data prevented the
required comparison with Wairakei drain-pot measurements, leaving
quantitative field agreement unresolved.
(\href{https://www.researchgate.net/publication/310818568_LIQUID_CARRYOVER_IN_GEOTHERMAL_STEAM-WATER_SEPARATORS}{Rizaldy
et al., 2016})

\begin{figure}[htbp]
\centering
\includegraphics[width=10cm]{\detokenize{/Users/shuheiyokkaichi/Developer/P4P_knowledgeWiki/Project/Final research report/Shuhei Report/figures/literature/rizaldy-2016-figure-9-entrainment.png}}
\caption*{Figure 5. Steam shear releases film droplets through (a) undercut and
(b) roll-wave entrainment. Source:
\href{https://www.researchgate.net/publication/310818568_LIQUID_CARRYOVER_IN_GEOTHERMAL_STEAM-WATER_SEPARATORS}{Rizaldy
et al.~(2016)}, fig.~9, crediting Ishii and Grolmes (1975).}
\end{figure}

Fluent roughness changes near-wall momentum transfer; it does not add a
film store or drain. EWF instead adds surface storage and transport.
Their outlet-flow responses must be read alongside retained mass and
collection.
(\href{https://ansyshelp.ansys.com/public/Views/Secured/corp/v252/en/flu_th/x1-5720008.164.html}{Ansys,
Inc., 2025k})

\subsection{2.5. Representing liquid and its
collection}\label{representing-liquid-and-its-collection}

The Mixture model represents bulk liquid through secondary-phase volume
fraction. Algebraic slip estimates relative phase velocity using local
equilibrium. Liquid volume fraction does not directly resolve film
thickness or droplet-size distribution.
(\href{https://ansyshelp.ansys.com/public/Views/Secured/corp/v252/en/flu_th/flu_th_sec_drift_oview.html}{Ansys,
Inc., 2025d})

DPM follows size-dependent droplet trajectories. Its interpretation
requires defined injection and boundary treatments, with every tracked
fate counted, including incomplete trajectories.

Skoog (2020) combined EWF surface transport, DPM deposition and
entrainment in a reactor-channel approximation. Film/droplet flows were
compared with one-dimensional calculations using experimentally based
Okawa correlations. Inlet allocation and transverse droplet velocity
affected the result. The separate liquid accounts provide a precedent;
reactor conditions and correlations still need justification for a
cyclone.
(\href{https://www.diva-portal.org/smash/get/diva2:1452100/FULLTEXT02.pdf}{Skoog,
2020};
\href{https://ansyshelp.ansys.com/public/views/secured/corp/v252/en/flu_ug/flu_ug_models_wallfilm.html}{Ansys,
Inc., 2025e})

Bulk-to-film transfer moves liquid within the domain; the absorber
removes it. Droplet injection must not duplicate another liquid feed.
Fluent permits cell-zone mass and associated transport sources, but
their law and effect require assessment. Section 3 describes the
contact-based absorber.
(\href{https://ansyshelp.ansys.com/public/Views/Secured/corp/v252/en/flu_ug/flu_ug_bcs_sec_cell_zones.html}{Ansys,
Inc., 2025a})

\subsection{2.6. Numerical startup, convergence and
validation}\label{numerical-startup-convergence-and-validation}

A steady calculation's iterations and pseudo-time describe numerical
development, not physical filling time. The National Aeronautics and
Space Administration (NASA) recommends monitoring residuals and relevant
outputs, with convergence criteria tied to acceptable output error.
(\href{https://www.grc.nasa.gov/www/wind/valid/tutorial/iterconv.html}{NASA,
2021a})

Fluent permits flow development before film initialisation. Reduced-feed
startup and ramp choices remain case-specific procedures requiring
assessment.
(\href{https://ansyshelp.ansys.com/public/views/secured/corp/v252/en/flu_ug/flu_ug_ewf_sec_overview.html}{Ansys,
Inc., 2025h})

Implementation checks, iterative convergence, mass accounting and
numerical sensitivity answer separate questions. Mesh and time-step
studies require a fixed model; roughness or absorber changes instead
test model sensitivity.
(\href{https://www.grc.nasa.gov/www/wind/valid/tutorial/verassess.html}{NASA,
2021c})

Physical validation requires relevant measurements and uncertainty
assessment. Here, field comparison must examine the assumed collection
route and outlet predictions within the tested operating range.
(\href{https://www.grc.nasa.gov/www/wind/valid/tutorial/valassess.html}{NASA,
2021b})

\subsection{2.7. Remaining model uncertainties and research
direction}\label{remaining-model-uncertainties-and-research-direction}

The literature makes the definition of collection central to the
research problem. A droplet that reaches a collecting boundary has a
terminal fate in Pointon's model, whereas the film literature allows
captured liquid to return to the steam. Purnanto's incomplete tracks
expose a further distinction between a calculated destination and an
assumed one. Agreement in an outlet measure therefore leaves the
internal liquid path unresolved when feed, capture or final fate is
imposed.

The inherited model adds a specific difficulty: its truncated domain
omits brine discharge. Mesh and inlet comparisons can change delivery to
the wall, but cannot supply the missing exit. This study therefore
connects those comparisons to collector development and then to film
storage, release and drainage. The separate bulk, film and droplet
accounts make it possible to ask whether a lower bulk-liquid outlet rate
accompanies external collection, internal transfer or continued storage.
That distinction determines what each result contributes towards
performance prediction.

\section{3. Model and methods}\label{model-and-methods}

\subsection{3.1. Computational domain and modelling
scope}\label{computational-domain-and-modelling-scope}

The model represented the upper separation region of a vertical cyclone:
spiral inlet, vessel walls and steam outlet. It omitted the brine pool
and discharge. The bottom remained a wall, so uncollected liquid could
accumulate or leave through the steam outlet.

A numerical collector represented transfer into the omitted pool without
opening a steam exit at the base. Full-geometry trials had not
established stable drainage and pool control (Appendix A.13). Pool
motion and downstream discharge remained outside the model.
Absorber-free cases disabled collection in the same domain.

\begin{figure}[htbp]
\centering
\includegraphics[width=13cm]{\detokenize{/Users/shuheiyokkaichi/Developer/P4P_knowledgeWiki/Project/Final research report/Shuhei Report/figures/model-domain-overview.png}}
\caption*{Figure 6. Solved 60,964-cell domain (a) and 6.754 mm wall-side inlet
strip (b). The lower collector is a numerical sink; the brine pool and
discharge are omitted.}
\end{figure}

The mesh uses vertical \(y\), with the bottom near \(y=0\) (exact bounds
in Appendix A.1). The collector contains 715 cells selected by centroid
position, \(y\leq0.10\) m; their nodes extend to about \(y=0.185\) m.
The full-feed parent, combined contact endpoint and shortened-startup
endpoint share coordinates and connectivity. The steam-outlet label
marks the computational boundary at the central pipe opening. Bulk flow
used a steady formulation, while film evolved on a separate clock. Bulk
startup described numerical development through iterations, not physical
startup time.

\subsection{3.2. Physical model and
assumptions}\label{physical-model-and-assumptions}

\subsubsection{3.2.1. Bulk steam--liquid
flow}\label{bulk-steamliquid-flow}

Ansys Fluent 2025 R2 used a pressure-based Mixture solver for bulk
continuity, momentum, liquid fraction and algebraic slip. Steam was
primary and water secondary. RNG \(k\)--\(\varepsilon\) included swirl
treatment and standard wall functions; gravity acted downwards. Energy
was disabled, excluding flashing, condensation and temperature-dependent
properties. Table 1 gives material values; Appendix A.1 records the
controls.
(\href{https://ansyshelp.ansys.com/public/Views/Secured/corp/v252/en/flu_th/flu_th_sec_multiphase_mixture.html}{Ansys,
Inc., 2025d})

\begin{table}[htbp]
\centering
\caption*{Table 1. Fixed material properties verified for the
throughput-controlled absorber model.}
{\setlength{\tabcolsep}{4pt}\begin{tabular}{@{}>{\raggedright\arraybackslash}p{4.13255cm}>{\raggedleft\arraybackslash}p{4.77255cm}>{\raggedleft\arraybackslash}p{6.53255cm}@{}}
\toprule
\textbf{Phase} & \textbf{Density (kg/m³)} & \textbf{Dynamic viscosity (Pa·s)} \\
\midrule
Vapour & 5.79743 & \(1.52062\times10^{-5}\) \\
Liquid water & 881.211 & \(1.45544\times10^{-4}\) \\
\bottomrule
\end{tabular}}
\end{table}

These three saved states used Manninen slip, Schiller--Naumann drag and
10 µm liquid-phase diameter. This local-equilibrium slip model was held
fixed during wall comparisons; its accuracy for concentrated wall liquid
remained unverified. The diameter controls continuum slip, separately
from DPM droplets and film (Appendix A.1).

\subsubsection{3.2.2. Wall roughness and wall-film
representation}\label{wall-roughness-and-wall-film-representation}

Roughness changed wall friction through equivalent sand-grain height
\(k_s\) and constant \(C_s\). Tested values were sensitivity inputs, not
measured surface properties. The wall treatment and unresolved settings
are documented in Appendix A.10.
(\href{https://ansyshelp.ansys.com/public/Views/Secured/corp/v252/en/flu_th/x1-5720008.164.html}{Ansys,
Inc., 2025k})

Eulerian Wall Film (EWF) solved surface mass and momentum, with
accretion transferring bulk liquid into the film. Through the 500 ms
reference, gravity, gas shear, wall viscosity and momentum advection
drove transport; native film-edge outflow supplied drainage. Additional
film forces and droplet interactions were disabled (Appendix A.9).
(\href{https://ansyshelp.ansys.com/public/Views/Secured/corp/v252/en/flu_ug/flu_ug_ewf_sec_options.html}{Ansys,
Inc., 2025e})

Later tests added pressure-gradient, spreading and surface-tension
forces, droplet collection, splashing, stripping and edge separation.
The last three mechanisms return film liquid to droplets. Changed bulk
fields prevent attribution to one mechanism. Film-to-bulk momentum
feedback stayed off, separately from numerical coupling of the film
equations (Appendix A.12).
(\href{https://ansyshelp.ansys.com/public/Views/Secured/corp/v252/en/flu_th/flu_th_ewf_sec_film_submodels.html}{Ansys,
Inc., 2025b})

\subsection{3.3. Boundary conditions and liquid
collection}\label{boundary-conditions-and-liquid-collection}

\subsubsection{3.3.1. Enthalpy and the inlet phase mass
flows}\label{enthalpy-and-the-inlet-phase-mass-flows}

The reference feed was 197.61 kg/s at 1600 kJ/kg and 11.2 bar absolute.
Saturated enthalpies gave vapour quality 0.408323. Inlet commands
retained Purnanto's rounded phase flows in Table 2; Appendix A.7 gives
the thermodynamic calculation.
(\href{https://www.researchgate.net/publication/269519626_CFD_MODELLING_OF_TWO-PHASE_FLOW_INSIDE_GEOTHERMAL_STEAM-WATER_SEPARATORS}{Purnanto
et al., 2013})

\subsubsection{3.3.2. Volumetric flow, inlet split and
velocity}\label{volumetric-flow-inlet-split-and-velocity}

The 0.724 m square inlet assigned about 0.933\% of its area to a 6.754
mm liquid strip, sized from phase volume flows under an assumed common
velocity. Appendix A.7 distinguishes the design properties from the
solved densities and resulting phase velocities.

The strip replaced an exploratory inlet that mixed liquid and steam over
half the opening. Upstream pipe films support testing liquid entry
beside the wall; complete segregation remained imposed. The matched
mixed/split comparison tested its consequences.
(\href{https://publications.mygeoenergynow.org/grc/1028587.pdf}{Pointon
et al., 2009})

\subsubsection{3.3.3. Fluent boundary
implementation}\label{fluent-boundary-implementation}

The wall-treatment calculations admitted pure liquid and vapour normal
to their respective inlets (Table 2). Operating pressure was 0 Pa, so
outlet gauge pressure equalled absolute pressure; inlet pressure was
calculated. Appendices A.6 and A.2 give turbulence/backflow inputs and
reconstruction differences.

\begin{table}[htbp]
\centering
\caption*{Table 2. Main boundary conditions. Reduced feed was 25\% of full feed,
with the phase mass ratio retained.}
{\setlength{\tabcolsep}{4pt}\begin{tabular}{@{}>{\raggedright\arraybackslash}p{3.30912cm}>{\raggedright\arraybackslash}p{4.26912cm}>{\raggedleft\arraybackslash}p{3.78912cm}>{\raggedleft\arraybackslash}p{3.78912cm}@{}}
\toprule
\textbf{Boundary} & \textbf{Condition} & \textbf{Full feed} & \textbf{Reduced feed} \\
\midrule
Liquid inlet & Liquid mass-flow inlet & 116.9200 kg/s & 29.2300 kg/s \\
Vapour inlet & Vapour mass-flow inlet & 80.6900 kg/s & 20.1725 kg/s \\
Steam outlet & Gauge pressure & 1.120 MPa & Same \\
Vessel and bottom & Stationary, no-slip walls & Selected roughness/film treatment & As defined for startup \\
\bottomrule
\end{tabular}}
\end{table}

\subsubsection{3.3.4. Contact-based
absorber}\label{contact-based-absorber}

The early absorber prescribed total removal from inlet throughput,
distributed by the liquid present in the collector. The later contact
absorber instead responded to local liquid inventory, without an
inlet-throughput cap or direct vapour mass sink:

\[
S_L=-\frac{\rho_L\alpha_L}{\tau},\qquad \tau=10^{-5}\ \mathrm{s}.
\]

Here, \(\rho_L\), \(\alpha_L\) and \(\tau\) are liquid density, volume
fraction and depletion time; \(S_L\) has units kg/(m³·s). Associated
momentum and turbulence sources represented one-way capture. Appendices
A.3--A.5 give the implementation, source checks, depletion-time
sensitivity and earlier throughput-controlled law.
(\href{https://ansyshelp.ansys.com/public/Views/Secured/corp/v252/en/flu_ug/flu_ug_bcs_sec_cell_zones.html}{Ansys,
Inc., 2025a})

Bulk removal, film drainage and diagnostic-particle escape were separate
routes. Early film drained where the active EWF wall ended. A later test
extended EWF onto the lower wall and added a direct film mass sink with
matching momentum removal and assumed 1.5 ms capture time. An isolated
source check and matched 15 ms OFF/ON runs tested it with all bulk
equations frozen (Appendix A.12).

\subsection{3.4. Mesh and spatial
discretisation}\label{mesh-and-spatial-discretisation}

Fluent Meshing produced polyhedral--hexahedral meshes with
wall-following layers for the expected liquid path. Reduced numerical
diffusion remained a hypothesis dependent on resolution, alignment and
discretisation.
(\href{https://ansyshelp.ansys.com/public/Views/Secured/corp/v252/en/flu_ug/flu_ug_GridTypes.html}{Ansys,
Inc., 2025c})

The main mesh contained 60,964 cells. Appendix A.6 gives spatial
schemes; Appendix A.10 gives connectivity and wall fields. Mean
main-wall \(y^+\) was approximately \(2.0\times10^4\)--\(3.4\times10^4\)
across three saved states, so wall-following topology did not establish
fine wall resolution. Layer dimensions and solver mesh-quality measures
remain unverified. Section 3.6 defines the separate mesh comparisons.

\subsection{3.5. Numerical solution
procedure}\label{numerical-solution-procedure}

Original startup used Hybrid Initialisation, smooth walls, EWF off and
throughput-controlled collection. A failed command left 25\% feed to
iteration 1580. The corrected 2000-update ramp preceded the change from
SIMPLE to Coupled with Global Time Step. Warm-up and a 1000-update
control supplied the iteration-5586 parent (Appendix A.5).

The \textbf{developed parent} was a saved full-feed field selected after
reviewing residuals, inventories, flows and sources. It supplied common
initial fields for wall comparisons, but retained a substantial
liquid-balance error and did not establish convergence. Shortened
startup reused iteration-1580 fields without reinitialisation and
enabled contact collection, 0.5 mm roughness, Coupled flow and dry film
together (Appendix A.16).

Coupled solved pressure and momentum corrections together, with Mixture
volume fraction solved separately. Coupling and pseudo-time changed
together, so their individual effects were not isolated. Later film runs
froze bulk flow, turbulence, volume fraction and relative motion.
Fluent's steady EWF procedure assumes a converged frozen carrier; here
the freeze allowed development under fixed forcing before bulk
convergence was established.
(\href{https://ansyshelp.ansys.com/public/Views/Secured/corp/v252/en/flu_ug/flu_ug_sec_uns_solve_pvel_1.html}{Ansys,
Inc., 2025i};
\href{https://ansyshelp.ansys.com/public/Views/Secured/corp/v252/en/flu_ug/flu_ug_sec_multiphase_solution.html}{2025j};
\href{https://ansyshelp.ansys.com/public/Views/Secured/corp/v252/en/flu_th/flu_th_ewf_sec_sol_alg.html}{2025g})

The implicit beta algorithm lacked inner residual reports. Equal-time
fields and mass accounts supplied local checks; missing residuals were
not passes. Save/reopen checked settings (Appendix A.9).

\subsection{3.6. Comparison design}\label{comparison-design}

Exploratory inlet, droplet and film tests guided later controlled
comparisons (Appendix A.13). Wall screens started from a common parent.
Startup and combined continuations changed several settings together and
tested procedures (Appendix A.16). Roughness ranged from 0 to 8 mm;
\(C_s\) variations and reporting windows are given in Appendix A.5.

The common parent held the inherited flow field fixed at the start of
each wall screen, reducing the influence of different preparation
histories. Equal update horizons and reporting windows then compared the
response to the selected treatment. Because the parent was still
developing, the comparison assessed that response rather than
equilibrium performance. The startup study asked a different question:
whether the combined model could reach full feed with smaller loading
disturbances. Its matched ramp progress aligned inlet loading, while the
separate activation hold retained the disturbance caused by introducing
the additional models.

Absorber-free mixed/split Coupled reconstructions held mesh, feed and
controls fixed, with fresh initialisation and 10,000 updates. The
historical tetrahedral comparison changed mesh, inlet and preparation.
SIMPLE/Coupled reconstruction comparisons also changed pseudo-time and
the \(k\) scheme, so neither isolated one cause.

Finer absorber-free SIMPLE diagnostics compared 342,609-, 679,970- and
997,604-cell fields at iteration 3000. These developing solutions were
not a converged mesh study (Appendix A.2).

The final mesh study compares five resolutions with the combined
wall-film and collection model, 0.045 mm roughness and film-to-flow
momentum feedback off (Section 4.13). The blocker and ramp address
separate geometry questions. Mesh independence remains unresolved.

\subsection{3.7. Data reduction and numerical
assessment}\label{data-reduction-and-numerical-assessment}

Each value retained its case, active equations and reporting window.
Endpoints, means and residual minima were distinct. Inventories, outlet
flows, collection and film-transfer rates were assessed with residuals
and source checks. A fixed iteration horizon did not establish
convergence. Appendix A.9 defines the numerical criteria and integration
methods.

Fluent fluxes were positive into the domain. With EWF off and no other
liquid transfer, the signed remainder (kg/s) was

\[
R_L=\sum_b\dot m_{L,b}+\dot m_{L,\mathrm{source,applied}}.
\]

The source retained its negative removal sign. Evaluated and applied
removal described the same sink, so only one entered the balance.
Bulk-to-film accretion cancelled only in compatible combined accounts.
Native mixture and phase-summed reports were assessed separately.

The early film balance compared integrated accretion with drainage plus
inventory change over accepted film time; Appendix A.9 defines the
integration and normalised error. Later release mechanisms required
separate transfer terms (Appendix A.12). With bulk frozen, reciprocal
bulk depletion did not advance, so film-account agreement could not
establish combined conservation.

Wall screens used final-500 arithmetic means; the 500 ms film run used
overlap-weighted final-10-ms rates. The bulk-liquid outlet ratio used
liquid feed as denominator; outlet dryness used total outlet flow
(Section 2.2). Neither pathlines nor diagnostic-droplet fates replaced
bulk-liquid fluxes. Residual scaling was checked for startup
comparisons; residuals were not percentage mass imbalances.
(\href{https://www.grc.nasa.gov/www/wind/valid/tutorial/iterconv.html}{NASA,
2021a})

\subsection{3.8. Experiment workflow and use of
automation}\label{experiment-workflow-and-use-of-automation}

The cycle was question → bounded comparison → implementation check →
Fluent solve → interpretation → next question. The researcher set
scientific scope and interpreted findings. Task-specific agent
instructions, Python and PyFluent supported execution. Native reports,
transcripts and reopened case/data pairs recorded what ran. This
workflow supported repeatable tests and recovery; numerical and physical
assessments remained separate. Appendix B links the records.

\section{4. Results}\label{results}

Mesh and inlet comparisons first assess liquid delivery and retention.
Collector and startup studies then supply fields for wall comparisons.
Film histories finally test whether captured liquid drains, remains
stored or returns to droplets. Earlier exploratory studies are retained
in Appendix A.13.

Each comparison uses its stated mesh and settings. Outlet liquid flow is
positive outwards unless a signed value is given. Steady iteration
describes numerical development; accepted film time describes the
transient film calculation.

\subsection{4.1. Mesh development and the wall-liquid
field}\label{mesh-development-and-the-wall-liquid-field}

The polyhedral--hexahedral mesh had larger interior cells, polyhedral
transitions and thin layers along the wall (Figure 7). The workflow gave
control over the mesh near the expected liquid path; Section 3.4 gives
its near-wall resolution.

At iteration 10,000, the historical tetrahedral field showed thin
wall-side liquid enrichment. The 60,964-cell SIMPLE reconstruction
showed a broader region with a higher liquid fraction. Inlet
representation and preparation also differed, and the coarse SIMPLE
state had a large mass-balance error. The comparison shows changed
fields; reduced numerical diffusion remains a possible explanation, not
a measured accuracy gain.

\begin{figure}[htbp]
\centering
\includegraphics[width=16cm]{\detokenize{/Users/shuheiyokkaichi/Developer/P4P_knowledgeWiki/Project/Final research report/Shuhei Report/figures/results-mesh-spatial-comparison.png}}
\caption*{Figure 7. Mesh sections (a--b) and liquid-fraction fields (c--d) at
iteration 10,000. The reconstructed case shows broader wall enrichment;
mesh, inlet and preparation differ, so the comparison does not isolate a
mesh effect.}
\end{figure}

\subsection{4.2. Two-phase inlet and wall-side liquid
placement}\label{two-phase-inlet-and-wall-side-liquid-placement}

Split feed produced a more continuous liquid-rich wall band and deeper
lower-region enrichment than mixed feed (Figure 8). At nominal 26.81
m/s, iteration-10,000 bulk-liquid inventory increased from 1262.6 to
1734.5 kg: a difference of 471.9 kg, or 37.4\%. Split feed retained more
liquid at all five tested speeds, 20.11--32.14 m/s.

Both cases used the same 60,964-cell mesh, total phase feed, Coupled
settings, initialisation and 10,000-update horizon. Mixed feed supplied
both phases across both inlet faces; split feed supplied liquid through
the wall-side strip. Complete phase segregation was an assumed profile
(Section 3.3). Over iterations 9500--10,000, outward steam-outlet liquid
averaged 99.671\% of the 116.9387 kg/s liquid feed for mixed feed and
99.404\% for split feed. The inventory difference therefore accompanied
almost complete liquid discharge through the steam outlet, rather than a
separate collection route.

\begin{figure}[htbp]
\centering
\includegraphics[width=16cm]{\detokenize{/Users/shuheiyokkaichi/Developer/P4P_knowledgeWiki/Project/Final research report/Shuhei Report/figures/results-inlet-spatial-comparison.png}}
\caption*{Figure 8. Wall-side split feed produces a more continuous liquid-rich
wall band than mixed feed. Matched iteration-10,000 sections use a
shared 0--1 scale; both cases have no collector or EWF.}
\end{figure}

\subsection{4.3. Liquid accumulation and the need for
collection}\label{liquid-accumulation-and-the-need-for-collection}

Saved fields from the three finer split-inlet SIMPLE cases placed liquid
near the wall and base. Their inventory and outlet histories show
accumulation with little discharge through the steam outlet (Figure 9).
At iteration 3000, the 342,609- and 679,970-cell inventories were both
about 454--455 kg, whereas the 997,604-cell case held 333.142 kg. The
latter discharged only 0.003623 kg/s through the steam outlet and held
69.16\% of its liquid below \(y=1\) m. These developing checkpoints do
not form a mesh-convergence result; Appendix A.17 retains their exact
values.

\begin{figure}[htbp]
\centering
\includegraphics[width=16cm]{\detokenize{/Users/shuheiyokkaichi/Developer/P4P_knowledgeWiki/Project/experiments/phase-08-storyline-reconstruction/stage-02-fine-mesh-simple/direct-student-342k/n3000/figures/matched-mesh-checkpoints.png}}
\caption*{Figure 9. Liquid accumulates while outlet flow remains small in the
finer absorber-free SIMPLE cases. Lines connect saved checkpoints;
solver iterations are not physical time.}
\end{figure}

On the 679,970-cell mesh, inventory reached 764.744 kg at iteration
4000, with 69.62\% below \(y=2\) m. Outlet liquid was 8.00 kg/s, or
6.84\% of the feed. The next continuation failed after printed iteration
4211; its cause was not isolated.

The inlet and mesh studies changed delivery and retention but supplied
no separate collection route. The next studies introduced a
liquid-selective source after unresolved full-geometry drainage trials
(Appendix A.13), then tested startup before comparing wall treatments.

\subsection{4.4. Reduced-feed development and inlet
ramp}\label{reduced-feed-development-and-inlet-ramp}

The startup and wall studies used the 60,964-cell split-inlet model,
with full feed of 116.92 kg/s liquid and 80.69 kg/s steam.

At 25\% feed with throughput collection, inventory fell only 0.30\% over
iterations 1080--1580, ending at 31.3571 kg (Figure 10). Minimum scaled
continuity was \(3.2056\times10^{-4}\).

A failed ramp command had left the model at reduced feed until iteration
1580. The corrected 2000-update ramp reached full feed at iteration
3580. Ramp continuity peaked at \(1.1166\times10^{-2}\) and inventory
reached 182.5178 kg.

\begin{figure}[htbp]
\centering
\includegraphics[width=12cm]{\detokenize{/Users/shuheiyokkaichi/Developer/P4P_knowledgeWiki/Project/Final research report/Shuhei Report/figures/results-described-early-development.png}}
\caption*{Figure 10. Reduced-feed hold, loading ramp and full-feed development. A:
ramp start; B: full feed and Coupled activation; C: EWF activation.
Residuals, inventory and inlet loading are shown against solver
iteration.}
\end{figure}

\subsection{4.5. Full-feed control and mass
balance}\label{full-feed-control-and-mass-balance}

The smooth-wall, EWF-off Coupled continuation became the iteration-5586
comparison parent. Continuity fell and inventory approached 296 kg, but
residual oscillations, outlet reverse flow and turbulent-viscosity
limiting persisted.

Applied throughput-controlled absorber removal matched its 116.92 kg/s
command, but another 24.3344 kg/s left through the steam outlet. The
signed liquid boundary-plus-source remainder was therefore −24.3344
kg/s, or −20.81\% of the feed. Source application was verified;
conservation remained unresolved (Appendix A.15.2).

This developing field supplied a common parent for the following wall
comparisons. Outlet-flow changes must be read with inventory and removal
because the parent did not establish a closed liquid account.

\subsection{4.6. Roughness response from the common
parent}\label{roughness-response-from-the-common-parent}

The roughness response was non-monotonic (Figure 11). At \(C_s=0.5\),
0.05 mm raised mean outlet flow by 21.21\%; 8 mm reduced it by 44.09\%,
from 24.3715 to 13.6253 kg/s. Increasing \(C_s\) lowered flow at both
fixed heights.

\begin{figure}[htbp]
\centering
\includegraphics[width=14cm]{\detokenize{/Users/shuheiyokkaichi/Developer/P4P_knowledgeWiki/Project/Final research report/Shuhei Report/figures/results-described-roughness.png}}
\caption*{Figure 11. Roughness height gives a non-monotonic outlet response; the 8
mm case also loses bulk inventory. EWF is off. Means use iterations
8087--8586; dashed lines mark the smooth-wall reference.}
\end{figure}

The 8 mm case lost 187.109 kg of bulk liquid. Oscillations and
incomplete balances prevented an efficiency claim. The later 0.045 mm
continuation also changed bulk advancement and film stepping, so it did
not isolate roughness (Appendix A.12).

\subsection{4.7. Accretion-enabled wall-film
response}\label{accretion-enabled-wall-film-response}

From common smooth-wall fields with throughput collection, the EWF
package reduced mean outlet liquid by 92.88\%, from 24.369 to 1.735 kg/s
over iterations 8087--8586 (Figure 12). Phase accretion and coupled film
equations were on; momentum feedback was off. Earlier trials are in
Appendix A.11.

\begin{figure}[htbp]
\centering
\includegraphics[width=14cm]{\detokenize{/Users/shuheiyokkaichi/Developer/P4P_knowledgeWiki/Project/Final research report/Shuhei Report/figures/results-described-film-outlet.png}}
\caption*{Figure 12. Accretion reduces bulk-liquid outlet flow; changing
inventories prevent a collection claim. Smooth-wall cases share the
full-feed parent. Shading marks the final 500 updates; the film step is
10 µs.}
\end{figure}

Bulk mass fell from 295.8536 to 63.0234 kg, while film mass reached only
3.1112 kg. The sections showed less liquid beside the wall (Figure 13).
The film package changed liquid distribution, but changing inventories
and an open combined liquid account prevent an efficiency claim.

\begin{figure}[htbp]
\centering
\includegraphics[width=14cm]{\detokenize{/Users/shuheiyokkaichi/Developer/P4P_knowledgeWiki/Project/Final research report/Shuhei Report/figures/bulk-comparison-described-film-treatment.png}}
\caption*{Figure 13. Less bulk liquid remains beside the wall with accretion
enabled. Both iteration-8586 centre sections use a shared 0--1 scale.
Film storage and drainage are assessed separately.}
\end{figure}

The following tests changed collection and startup, then extended film
development to examine drainage and storage.

\subsection{4.8. Contact collection and combined model
development}\label{contact-collection-and-combined-model-development}

Shorter depletion time reduced collector inventory but increased removal
variability (Appendix A.4). The combined continuation applied 10 µs
contact collection, 0.5 mm roughness and a 1 µs film step together;
their effects were not isolated.

At iteration 45606, film mass was 6.3600 kg after 120.769 ms.
Final-1000-update drainage remained 9.98\% below accretion, so the film
continued to fill. Appendix A.15.3 and the linked case history retain
the endpoints and detailed film account.

The shortened-startup comparison below tests whether introducing these
treatments before loading reduces the ramp disturbance.

\subsection{4.9. Shortened startup with early wall-film
activation}\label{shortened-startup-with-early-wall-film-activation}

Shortened startup introduced the combined wall and collection treatment
before loading. Relative to the historical ramp, peak continuity fell by
70.15\%, peak bulk + film inventory by 70.43\%, and peak outward outlet
liquid by 75.78\% (Figure 14). Appendix A.17 retains the peak values.

\begin{figure}[htbp]
\centering
\includegraphics[width=16cm]{\detokenize{/Users/shuheiyokkaichi/Developer/P4P_knowledgeWiki/Project/Final research report/Shuhei Report/figures/results-described-startup.png}}
\caption*{Figure 14. Shortened startup reduces ramp peaks. (a) Continuity; (b)
bulk + film inventory; (c) signed outlet liquid flow, negative outwards.
The larger activation disturbance is excluded; dashed line: full feed;
shading: subsequent hold.}
\end{figure}

The excluded activation hold had a continuity peak of 0.10593, exceeding
the historical ramp peak. At iteration 5080, bulk and film masses were
61.0549 and 0.1645 kg. Almost all film input was still stored (Appendix
A.15.4). All ramp updates met the recorded inner-film tolerance, but 15
of the final 100 hold updates failed it. The recipe reduced loading
peaks without establishing complete startup convergence (Appendix
A.9.4).

The film was still young at full feed. The following continuation held
the bulk fields fixed to examine its development under unchanged
forcing.

\subsection{4.10. Film development under frozen bulk
fields}\label{film-development-under-frozen-bulk-fields}

Film mass increased from 0.1645 kg at 3.5 ms to 9.5444 kg at 500 ms
(Figure 15). After iteration 5080, bulk equations were frozen, leaving
bulk inventory at 61.0549 kg and outward bulk-liquid outlet flow at
3.4300 kg/s. These fixed values describe held forcing; reciprocal bulk
depletion was not calculated.

\begin{figure}[htbp]
\centering
\includegraphics[width=16cm]{\detokenize{/Users/shuheiyokkaichi/Developer/P4P_knowledgeWiki/Project/Final research report/Shuhei Report/figures/results-film-development-reference500ms.png}}
\caption*{Figure 15. Drainage develops while film inventory keeps rising under
frozen bulk forcing. (a) Film mass; (b) input and drainage. Shading
marks the final 10 ms; markers locate the Figure 16 snapshots and
endpoint.}
\end{figure}

The film spread and thickened over the lower and middle wall (Figure
16). Between 75.66 and 265.18 ms, wall area with thickness at least 0.10
mm increased from 40.34\% to 61.98\%; maximum thickness rose from 0.253
to 0.291 mm.

\begin{figure}[htbp]
\centering
\includegraphics[width=16cm]{\detokenize{/Users/shuheiyokkaichi/Developer/P4P_knowledgeWiki/Project/Final research report/Shuhei Report/figures/wall-film-described-development.png}}
\caption*{Figure 16. Film spreads and thickens over the lower and middle wall.
Saved times and masses are labelled; all three views use the same camera
and 0--0.30 mm scale, with bulk fields frozen.}
\end{figure}

Over the final 10 ms at 500 ms, mean accretion was 81.121 kg/s and
drainage was 69.600 kg/s, a 14.20\% drainage deficit. Film storage,
calculated from mass change divided by accepted film time, remained
11.528 kg/s. The late mass curve therefore continued to rise despite
increasing drainage. Rates use interval-overlap weighting; Appendix A.17
retains the earlier 265.184 ms result and its different 1000-update
window. Only the reference-speed arm reached the 500 ms endpoint
(Appendix A.9).

Continued storage motivated tests of two further destinations: release
back to droplets and direct removal from the lower wall.

\subsection{4.11. Added film-release mechanisms and direct
drainage}\label{added-film-release-mechanisms-and-direct-drainage}

A separate configuration enabled stripping, splashing, droplet
collection and edge separation, retaining film-to-flow momentum feedback
off. Its bulk fields and film forces differed from the 500 ms reference.
After 4000 bulk updates, 50 ms of frozen-bulk development increased film
mass from 0.2638 to 5.9070 kg. Final approximately 5 ms mean stripping
and edge-separation rates were 64.56 and 51.87 kg/s, while storage was
80.81 kg/s. A separate splash contribution was unresolved.

Native film outflow was zero during that interval. Extending the film
domain onto the lower wall then allowed a matched direct-sink test
(Figure 17). Drain OFF and ON started from identical fields and 9.7526
kg total film, with the same expanded film domain and all bulk equations
frozen. Over 15 ms, the sink reduced final lower-wall film by 70.32\%.
Native outflow also fell, so the 0.5088 kg direct removal increased
combined external removal by only 0.1897 kg.

\begin{figure}[htbp]
\centering
\includegraphics[width=14cm]{\detokenize{/Users/shuheiyokkaichi/Developer/P4P_knowledgeWiki/Project/Final research report/Shuhei Report/figures/results-direct-film-drain.png}}
\caption*{Figure 17. Direct drainage reduces lower-wall film but partly replaces
native outflow. Matched 15 ms OFF/ON test with bulk frozen: (a)
lower-wall mass; (b) native outflow and direct removal.}
\end{figure}

In the longer continuation, lower film stayed near 0.056 kg while upper
film grew. The complete 405 ms record, including the rejected tail,
contains 15.16 kg of direct removal and a 27.58 kg increase in film
inventory. Storage was positive before the Courant-guard crossing. Local
drainage therefore did not stop accumulation elsewhere. Particle release
was internal transfer, while the sink removed liquid from the model
(Appendix A.12).

The following numerical checks assess source application and film-step
selection. The final mesh comparison addresses resolution separately.

\subsection{4.12. Numerical assessment of the selected development
path}\label{numerical-assessment-of-the-selected-development-path}

The shortened-startup contact source matched independently evaluated
removal at all 3500 solved updates. A separate film-step test from
identical fields showed that 20 µs passed local field and ledger checks
over 2.5 ms, whereas 25 µs exceeded the 0.1\% ledger-error limit.

These checks supported implementation and local step selection.
Inner-film failures, unavailable residuals and continued storage
prevented a convergence claim (Appendix A.9.4).

\clearpage
\proofpage

\subsection{4.13. Final-model mesh
sensitivity}\label{final-model-mesh-sensitivity}

The combined wall-film and collection model is being assessed on meshes
containing 60,964, 342,609, 679,970, 997,604 and 2,596,657 fluid cells.
The comparison asks whether wall-liquid distribution, film storage and
removal change with resolution. A completed comparison at matched film
ages and bulk-development states is not yet available; mesh independence
remains unresolved.

\begin{quote}
\textbf{Figure 18 reserved --- one combined mesh comparison.} Present
the most informative bulk-liquid or film-thickness views at shared
scales, with a compact panel for the corresponding storage and removal
response. Use common geometry, loading, wall treatment and reporting
windows. Identify any unmatched or incomplete cases.
\end{quote}

\proofslot{9.0cm}{Figure 18: final-model mesh sensitivity.\\Matched liquid fields with a compact storage/removal comparison.\\Results pending.}

The final finding will distinguish spatial changes from continued film
development and quantify their effect on liquid storage, collection and
steam-outlet flow. Detailed case settings and balance terms belong in
the supporting records. The vertical blocker and angled ramp change
geometry and are assessed separately.

\vfill\noindent\textit{Review layout only. Replace this allowance with the verified mesh comparison and its finding.}

\clearpage

\section{5. Discussion}\label{discussion}

The results identify three linked model sensitivities: delivery to the
wall, transfer into film and removal from the vessel. The discussion
connects those responses to the collection question and examines the
explanations that remain unresolved.

\subsection{5.1. Mesh development changed the wall-liquid
representation}\label{mesh-development-changed-the-wall-liquid-representation}

The polyhedral--hexahedral workflow made wall refinement and coarse
local testing more practical. This matters because the liquid path of
interest occupies the wall region, where fine cross-wall spacing is
needed to represent steep liquid-fraction gradients. The wall layers
gave control over that spacing, with larger interior cells supporting
less costly exploratory calculations. Reduced numerical smearing remains
a hypothesis because transport error also depends on resolution,
alignment and discretisation.
(\href{https://ansyshelp.ansys.com/public/Views/Secured/corp/v252/en/flu_ug/flu_ug_GridTypes.html}{Ansys,
Inc., 2025c, §7.1.3.3})

The broader liquid-rich region in Figure 7 suggested that wall liquid
was being represented differently. The finer absorber-free fields then
placed much of the retained liquid near the wall and base while sending
little through the steam outlet (Figure 9). These observations support
the concern that a model which retains liquid away from steam also needs
a route for its collection. Their later failure did not establish
greater mesh accuracy or identify the cause of numerical breakdown.

Historical comparisons also changed inlet representation and
preparation, and lacked matched inventory histories. The combined-model
mesh study will test whether the wall-liquid fields, storage and removal
rates persist across resolutions under one common model (Section 4.13).
Residual reduction alone would not answer that question.

\subsection{5.2. Inlet phase placement changes the wall-liquid starting
condition}\label{inlet-phase-placement-changes-the-wall-liquid-starting-condition}

The thin liquid strip made the wall-rich feed assumption reproducible.
Films in upstream piping provide a physical reason to test that
placement, but do not establish the actual inlet profile.
(\href{https://publications.mygeoenergynow.org/grc/1028587.pdf}{Pointon
et al., 2009})

With feed, mesh and controls fixed, wall-side placement increased
retained liquid by 37.4\%. Figure 8 also showed a more continuous
liquid-rich wall band and deeper lower-region enrichment. This
establishes that inlet placement changes the liquid distribution
presented to the wall before EWF is added. Almost all supplied liquid
still left through the steam outlet, so greater retention did not
establish better collection. Later wall comparisons must preserve inlet
placement to separate wall-treatment effects from a changed starting
distribution, then follow that liquid through storage and removal.

\subsection{5.3. The central development problem was the liquid
collection
route}\label{the-central-development-problem-was-the-liquid-collection-route}

The closed-bottom model's only exit was the steam outlet. Its liquid
either accumulated or left with steam, whereas separator operation
requires outward transport followed by brine discharge. Rising inventory
was consistent with the missing drain, but changes over steady
iterations did not give a physical filling rate.
(\href{https://doi.org/10.1016/j.geothermics.2014.05.009}{Zarrouk \&
Purnanto, 2015})

The absorber-free Coupled cases discharged 99.4--99.7\% of liquid feed
through the steam outlet. Their smaller mass imbalance accompanied an
unsuitable liquid exit route. In a stationary closed-bottom calculation
without another sink, conservation would require incoming liquid to
leave through that outlet. Closure can therefore be numerically
desirable while the available boundary routes remain inadequate for the
separator question. This distinction explains why Coupled remained
useful after a collector was added.

Opening the lower boundary let steam leave below. Full-geometry trials
exposed dependence on a liquid seal and outlet control; neither an
initially patched pool nor the tested pressure feedback established a
maintained pool. Explicit collection in the truncated vessel allowed
wall studies to continue. It represents transfer to an omitted pool
without predicting its pressure, level or discharge response.

\subsection{5.4. Contact removal and staged startup made the model
easier to
develop}\label{contact-removal-and-staged-startup-made-the-model-easier-to-develop}

The contact absorber tied removal of mass and momentum to liquid present
in the collector. This better matched transfer to the omitted pool than
prescribing removal from inlet throughput alone.

Reduced feed allowed flow and collection to develop before full loading.
The useful hold arose after a ramp command failed; the checked ramp
turned that observation into a repeatable procedure. Activating the
later contact and film treatments before loading reduced ramp peaks in
continuity residual, combined bulk + film inventory and outlet flow by
about 70\%, 70\% and 76\%.

The shorter procedure supported finer-mesh repeats, but changed several
settings together. The larger activation disturbance occurred before the
ramp comparison, so the smaller ramp peaks do not describe the whole
startup. The causes of that disturbance and later film-solver failures
remain unresolved. Comparable starts require checked fields and film
ages.

\subsection{5.5. Wall treatment strongly controls the calculated liquid
state}\label{wall-treatment-strongly-controls-the-calculated-liquid-state}

The non-monotonic roughness response in Figure 11 shows that increasing
the roughness input did not consistently reduce outlet liquid. Small
heights raised outlet liquid flow, while the largest tested height
reduced it by 44.09\%. Roughness changes the modelled near-wall momentum
transfer; it does not supply a new liquid store or drain.
(\href{https://ansyshelp.ansys.com/public/Views/Secured/corp/v252/en/flu_th/x1-5720008.164.html}{Ansys,
Inc., 2025k}) One possible explanation is changed liquid transport or
residence near the wall. A competing explanation is reduced liquid
availability after bulk inventory depletion. The 8 mm case lost 187.109
kg during the continuation, so the latter explanation cannot be
dismissed because the final outlet history became comparatively level.
The wall-surface velocity report sampled the no-slip boundary and could
not resolve the near-wall transport needed to distinguish these
explanations.

The accretion-enabled film package changed the representation itself by
adding surface storage and transport. Mean bulk-liquid outlet flow fell
from 24.369 to 1.735 kg/s, and Figure 13 showed less bulk liquid beside
the wall. However, bulk mass fell from 295.854 to 63.023 kg while film
mass reached only 3.111 kg. The lower outlet flux and the new film
inventory therefore do not account for the liquid destination by
themselves. The unresolved bulk, source and film account prevents an
efficiency interpretation.

These screens identify wall treatment as a major model sensitivity. A
useful follow-up must distinguish increased delivery to the collector
from depletion, surface storage and transfer back to droplets. It needs
wall-adjacent fields and separate inventories and removal histories,
rather than an outlet ranking alone. The 8 mm screen and later 0.045 mm
material-reference continuation also used different bulk advancement and
film stepping. Neither identified an optimum surface finish.

\subsection{5.6. Film development, release and drainage answer different
parts of the same
question}\label{film-development-release-and-drainage-answer-different-parts-of-the-same-question}

Longer histories answered the earlier question about patchy initial
films by showing continued spreading and thickening. Between 75.66 and
265.18 ms, appreciable film coverage increased as well as maximum
thickness (Figure 16). By 500 ms, drainage had developed to 69.60 kg/s,
but input remained 81.12 kg/s and storage was still 11.53 kg/s over the
final 10 ms (Figure 15). A wider visible film and substantial drainage
can therefore coexist with continuing accumulation.

The bulk fields were frozen during this development, so the film
responded to held forcing. Constant bulk inventory and outlet flux were
imposed by that choice; reciprocal bulk depletion was not advanced. The
run reveals how the film develops under those fields, while sustained
coupled collection remains a separate question.

Later runs used different fields and film forces. Stripping and
edge-separation records showed film liquid returning to droplets.
Splashing was enabled but not measured separately. Release reduces film
mass without necessarily removing liquid from the vessel: droplets may
return to the wall, enter the collector or leave with steam. This
reflects the film-entrainment problem in the literature and explains why
wall capture alone cannot establish permanent collection.
(\href{https://www.researchgate.net/publication/310818568_LIQUID_CARRYOVER_IN_GEOTHERMAL_STEAM-WATER_SEPARATORS}{Rizaldy
et al., 2016};
\href{https://ansyshelp.ansys.com/public/Views/Secured/corp/v252/en/flu_th/flu_th_ewf_sec_film_submodels.html}{Ansys,
Inc., 2025b})

The matched direct-drain test proved local film removal with bulk
equations frozen (Figure 17). Lower-wall inventory fell by 70.32\%, but
the 0.5088 kg direct sink partly replaced native film outflow, which
fell from 0.5302 to 0.2111 kg. Combined removal increased by only 0.1897
kg, close to the reduction in total film inventory. The sink changed the
division between removal routes as well as adding removal. Its rate
alone would therefore overstate the gain relative to the drain-OFF arm.

In the longer continuation, lower inventory stayed near 0.056 kg while
almost all film growth occurred above it. Removal was active, yet input
still exceeded the combined losses; upper storage was positive before
the later Courant-guard failure. This supports a transport question as
well as a local-drain question. It does not establish that downward
motion was absent or isolate the cause of upper accumulation. Features
that intercept circulating film or guide it towards the collector are
specific design hypotheses. Their assessment must follow collection,
steam-outlet liquid, storage and pressure loss together, because lower
film mass could also result from greater release into steam.

\subsection{5.7. Contribution and the evidence needed for performance
prediction}\label{contribution-and-the-evidence-needed-for-performance-prediction}

The contribution is a model-development approach that follows liquid
through inlet placement, wall transport and collection. Checked
collection and startup procedures enabled the comparisons; longer
histories distinguished storage, release and drainage. Appendix B
summarises the supported answers in Table B2.

Bulk-liquid flux, outlet quality and DPM escape fractions measure
different quantities. Exact absorber tracking did not close the early
balance, and a small film-ledger error did not establish whole-separator
conservation. Each transfer must be followed and counted once.

High saved wall \(y^+\) values, incomplete refinement assessment and
omitted pool response limit performance prediction. Physical validation
still requires matched operating conditions, measurement locations and
uncertainty.
(\href{https://www.grc.nasa.gov/www/wind/valid/tutorial/valassess.html}{NASA,
2021b})

\section{6. Conclusions and future
work}\label{conclusions-and-future-work}

\subsection{6.1. Conclusions}\label{conclusions}

Liquid reaching the wall did not by itself establish collection. This
study developed a geothermal separator CFD model by following where
liquid enters, accumulates, transfers and leaves. Mesh and inlet
investigations showed that the predicted wall-liquid distribution
depended on its representation. In the matched inlet comparison, placing
liquid beside the wall increased retained mass by 37.4\%. Absorber-free
cases exposed the missing collection route through accumulation or
almost complete liquid discharge through the steam outlet.

A liquid-selective absorber and staged loading supplied common fields
for wall studies. The contact treatment represented transfer of liquid
mass and momentum to the omitted brine pool. The later combined startup
reduced the peak continuity residual during the ramp by 70.15\%, while
retaining a separate activation disturbance.

Wall treatment strongly changed the predicted liquid state. The selected
accretion-enabled film package reduced mean bulk-liquid outlet flow from
24.369 to 1.735 kg/s, alongside a large inventory change. Film
calculations showed spreading, thickening and drainage, with storage
still averaging 11.53 kg/s over the final 10 ms at 500 ms. The direct
drain removed 0.5088 kg but increased combined external removal by only
0.1897 kg because it partly replaced native film outflow. Its removal
alone therefore overstated the additional collection. Lower-wall
collection could operate while film continued to accumulate above it.

The completed work establishes model sensitivity and checked collection
and startup procedures. It also gives a specific reason to test designs
that redirect wall liquid towards collection. Mesh sensitivity of the
combined model remains unresolved. Sustained liquid accounting and
physical validation are still needed before claiming improved separation
efficiency or pressure-drop accuracy.

\subsection{6.2. Future work}\label{future-work}

The next comparisons should test the explanations established here:

\begin{enumerate}
\def\labelenumi{\arabic{enumi}.}
\tightlist
\item
  \textbf{Complete the mesh-sensitivity assessment.} Compare liquid
  distributions and rates using fixed geometry and model settings at
  matched film ages and bulk-development states. Follow up quantities
  that remain resolution-dependent.
\item
  \textbf{Complete the coupled liquid account.} Advance the bulk
  equations and track bulk, film and droplet transfers together.
  Separate external removal from internal transfers and repeat
  developed-film step checks.
\item
  \textbf{Separate startup and model effects.} Compare loading sequence,
  collector law and solver controls individually from common fields.
\item
  \textbf{Compare with measurements.} Match feed conditions and pressure
  and liquid-sampling locations to assess the model's use for design.
\end{enumerate}

Initial design work now applies the wall-film findings to two features.
A \textbf{vertical film blocker} is intended to intercept circulating
liquid before it meets incoming steam. An \textbf{angled film ramp} is
intended to guide that liquid towards lower collection. Both
calculations have reached 150 ms film age; their final 100 ms used
matching basic wall-film settings and 10 µs steps with bulk fields held
fixed (Appendix A.14). Their meshes, initial flow fields and earlier
film histories differ, so the results do not yet isolate a geometry
effect or establish improved separation. The next comparison should
follow collection, steam-outlet liquid, film storage and pressure loss
from controlled baselines. Including film-release mechanisms will also
test whether interception reduces renewed entrainment. Lower film
inventory alone cannot distinguish collection from release into steam.

\end{document}
